\section{Conclusion and future work}
\label{sec:conclusion}

This paper turned TrajRAG from an aviation method into a domain-agnostic,
training-free framework for reconstructing trajectories after
satellite-navigation spoofing has been detected. The method is split into a
shared engine (phase-aware retrieval, a frozen large language model that
proposes the next position, a kinematic feasibility guard and a feedback
loop) and a domain adapter made of a prompt template, phase-indexed
memories, a trusted non-satellite signal and constants obtained by a fixed
calibration procedure. Moving from aircraft to ships then required no model
training. On commercial vessels in two port corridors, evaluated against a
retrieval memory that structurally excludes both the test set and any
chronologically later voyage, the maritime instantiation reduced error by a
factor of 1.3--1.5 relative to the best classical baseline in every phase.
It was significantly more accurate than classical, recurrent and
retrieval-only methods in every phase tested. Against a trained,
ship-specific convolutional network, no significant difference was detected
in any phase (Houston, $n=50$); we read matching a network purpose-trained
on this corridor's own traffic, using no maritime training data at all, as
the relevant finding. TrajRAG reconstructs a full blackout faster than a
single AIS reporting interval. Together with the aviation results, this
supports a single reconstruction service for ATC and VTS that is extended
by adding adapters rather than by training models, with the caveat that its
accuracy advantage over a dedicated, well-resourced domain model is not yet
established on this evidence.

Future work will (i) evaluate the phase-specific tolerances already derived
by the calibration procedure ($\alpha \approx 1.5$--2.0 across phases on the
corrected data), which requires new reconstructions rather than re-analysis
of existing results, since the guard sits inside the autoregressive loop
(Section~\ref{sec:maritime-calibration}); (ii) re-run the backbone-sensitivity
and ablation-forecaster comparisons at full scale under the corrected
preprocessing and retrieval memory, which budget limited this paper to a
smaller sample; (iii) extend the forecasters to the full horizon and
retrain them, and AIS-ACNet, on the corrected data at the full 125-voyage
scale to recover the statistical power this paper's budget could not fund;
(iv) validate on documented maritime interference incidents with
independent radar or satellite references; (v) instantiate further domains
to test how far the adapter abstraction extends; and (vi) implement and
evaluate online, trusted-signal-only phase identification for maritime,
rather than relying on offline full-voyage segmentation
(Section~\ref{sec:threats}).

\section*{CRediT authorship contribution statement}
\textbf{Yuval Felendler:} Conceptualization, Methodology, Software,
Investigation, Data curation, Formal analysis, Visualization, Writing --
original draft. \textbf{Ruben Sasson:} Software, Investigation, Validation,
Writing -- review \& editing. \textbf{Yuval Elovici:} Supervision,
Resources, Writing -- review \& editing. \textbf{Asaf Shabtai:}
Conceptualization, Supervision, Methodology, Writing -- review \& editing.

\section*{Declaration of competing interest}
The authors declare that they have no known competing financial interests
or personal relationships that could have appeared to influence the work
reported in this paper.

\section*{Data availability}
The maritime corpus is built from public NOAA Marine Cadastre AIS archives
\citep{noaa2025marinecadastre} and Open-Meteo Marine data
\citep{openmeteo2025}. Code for both instantiations, all baselines, the
calibration analysis and the statistical tests, together with the voyage
identifiers of the test sets, is available at
\url{https://github.com/yf2595/TrajRAG}. Aviation trajectories are subject
to FlightRadar24 and Meteomatics terms; the exact queries needed to rebuild
the corpus are documented in the repository. LLM-based experiments require
access to the OpenAI models named in Section~\ref{sec:implementation}.

\section*{Declaration of generative AI and AI-assisted technologies in the
manuscript preparation process}
During the preparation of this work the authors used an AI writing
assistant to edit and restructure the manuscript text. After using this
tool, the authors reviewed and edited the content as needed and take full
responsibility for the content of the published article. Large language
models used as components of the proposed method are described in
Section~\ref{sec:implementation}.

\section*{Funding}
This research did not receive any specific grant from funding agencies in
the public, commercial, or not-for-profit sectors.
